\documentclass[a4paper,12pt,twoside]{memoir}
\sloppy
% Castellano
\usepackage[spanish,es-tabla]{babel}
\selectlanguage{spanish}
\usepackage[utf8]{inputenc}
\usepackage[T1]{fontenc}
\usepackage{lmodern} % Scalable font
\usepackage{microtype}
\usepackage{placeins}
% Añado paquete para poder añadir TODO's
\usepackage{todonotes}
\usepackage{tabularx}
\usepackage{makecell}

\RequirePackage{booktabs}
\RequirePackage[table]{xcolor}
\RequirePackage{xtab}
\RequirePackage{multirow}

% Añado listings para ver el código en la memoria mejor
\usepackage{listings}

% Definir colores personalizados
\lstset{
	basicstyle=\ttfamily\footnotesize,  % Estilo básico del texto
	keywordstyle=\color{blue},          % Estilo de las palabras clave
	commentstyle=\color{green},         % Estilo de los comentarios
	stringstyle=\color{red},            % Estilo de las cadenas
	numbers=left,                       % Números de línea en la izquierda
	numberstyle=\tiny\color{gray},      % Estilo de los números de línea
	stepnumber=1,                       % Cada línea numerada
	frame=single,                       % Cuadro alrededor del código
	breaklines=true,                    % Romper líneas largas
	captionpos=b,                       % Posición del caption (debajo)
	tabsize=2,                          % Tamaño de tabulación
	showspaces=false,                   % No mostrar espacios
	showstringspaces=false,             % No mostrar espacios en las cadenas
	showtabs=false,                     % No mostrar tabulaciones
	language=Python		                % Se usa el estilo de Python como base
}

\lstdefinelanguage{JavaScript}{
	keywords={break, case, catch, continue, debugger, default, delete, do, else, finally, for, function, if, in, instanceof, new, return, switch, this, throw, try, typeof, var, let, const, while, with},
	keywordstyle=\color{purple}\bfseries,
	ndkeywords={class, export, boolean, throw, implements, import, this},
	ndkeywordstyle=\color{blue}\bfseries,
	identifierstyle=\color{black},
	sensitive=true,
	comment=[l]{//},
	morecomment=[s]{/*}{*/},
	commentstyle=\color{gray}\ttfamily,
	stringstyle=\color{orange}\ttfamily,
	morestring=[b]',
	morestring=[b]"
}

\lstdefinestyle{jsStyle}{
	language=JavaScript,
	basicstyle=\ttfamily\footnotesize,
	keywordstyle=\color{purple},
	commentstyle=\color{gray},
	stringstyle=\color{orange},
	numbers=left,
	numberstyle=\tiny\color{gray},
	stepnumber=1,
	frame=single,
	breaklines=true,
	captionpos=b,
	tabsize=2,
	showspaces=false,
	showstringspaces=false,
	showtabs=false
}


\renewcommand\lstlistlistingname{Índice de Códigos}
\renewcommand\lstlistingname{Código}



% Links
\PassOptionsToPackage{hyphens}{url}\usepackage[colorlinks]{hyperref}
\hypersetup{
	allcolors = {red}
}

% Acrónimos y glosario
\usepackage[acronym]{glossaries}
\makenoidxglossaries
%\usepackage[acronym]{glossaries}
%\makeglossaries

\newacronym{APEX}{APEX}{\textit{Application Express}}
\newacronym{TFG}{TFG}{\textit{Trabajo de Fin de Grado}}
\newacronym{ODB}{ODB}{\textit{ORACLE 23ai}}
\newacronym{FN}{FN}{\textit{Forma Normal}}
\newacronym{SGBD}{SGBD}{\textit{Sistemas de Gestión de BD}}
\newacronym{HTTPS}{HTTPS}{\textit{Secure Hipertext Transfer Protocol}}
\newacronym{VYP}{VYP}{\textit{Validación y Pruebas}}
\newacronym{BDR}{BDR}{\textit{Base de Datos Relacional}}
\newacronym{PK}{PK}{\textit{Primary Key-Clave Primaria}}
\newacronym{FK}{FK}{\textit{Foreing Key-Clave Externa}}
\newacronym{IDE}{IDE}{\textit{Entorno de Desarrollo Integrado}}
\newacronym{DDL}{DDL}{\textit{Lenguaje de Definición de Datos}}
\newacronym{DML}{DML}{\textit{Lenguaje de Manipulación de Datos}}
\newacronym{SQL}{SQL}{\textit{Structured Query Language}}
\newacronym{MPV}{MPV}{\textit{Mínimo Producto Viable}}
\newacronym{ODS}{ODS}{\textit{Objetivos de Desarrollo Sostenible}}
\newacronym{ORDS}{ORDS}{Oracle REST Data Services}
\newacronym{REST}{REST}{REpresentational State Transfer}
\newacronym{SSL}{SSL}{Secure Sockets Layer}
\newacronym{OCI}{OCI}{Oracle Cloud Infrastructure}
\newacronym{VM}{VM}{Máquina Virtual}
\newacronym{AOP}{AOP}{APEX Office Print}
\newacronym{PL/SQL}{PL/SQL}{Procedural Language/Structured Query Language}
\newacronym{IA}{IA}{Inteligencia Artificial}
\newacronym{ACL}{ACL}{Listas de Control de Acceso}
\newacronym{JS}{JS}{JavaScript}
\newacronym{UXXI}{UXXI}{Universitas XXI}

% Glosario
\newglossaryentry{Milestone}{
	name=Milestone,
	description={Punto de control o evento significativo dentro de un proyecto que marca la finalización de una fase importante o la consecución de un objetivo específico}
}

\newglossaryentry{Copilot}{
	name=Copilot,
	description={Asistente de inteligencia artificial desarrollado por Microsoft, diseñado para ayudar a los usuarios en diversas tareas a través de diferentes aplicaciones}
}

\newglossaryentry{Datos MOCK}
{
	name=Datos MOCK,
	description={Datos ficticios para probar aplicaciones, bases de datos, interfaces, sin necesidad de usar datos reales o confidenciales}
}

\newglossaryentry{LATEX}
{
	name=LATEX,
	description={Sistema de composición de textos orientado a la creación de documentos escritos, genralmente científicos que presenten una alta calidad tipográfica}
}

\newglossaryentry{Cardinalidad}
{
	name=Cardinalidad,
	description={Relación entre el número de instancias de una entidad con respecto a otra entidad con la que se relaciona}
}

\newglossaryentry{CasodeUso}
{
	name=Caso de Uso,
	description={Descripción de cómo los actores(usuarios), interactúan con el sistema para lograr un objetivo específico}
}

\newglossaryentry{Atomicidad}
{
	name=Atomicidad,
	description={Propiedad de un campo, por el cual no puede ser descompuesto en pedazos más pequeños por el DBMS}
}

\newglossaryentry{Normalización}
{
	name=Normalización,
	description={Proceso de estructuración los datos en una BD relacional, eliminando redundancia de datos y manteniendo la integridad a través de Formas Normales}
}

\newglossaryentry{MER}
{
	name=Modelo Relacional,
	description={Herramienta que permite organizar los datos en tablas y las relaciones entre las mismas}
}

\newglossaryentry{SLL}
{
	name=Secure Sockets Layer,
	description={protocolo de seguridad que cifra la comunicación entre un navegador web y un servidor, protegiendo la información transmitida}
}

\newglossaryentry{Release}
{
	name= Release,
	description={Versión específica de un software que ha sido completada, probada y puesta a disposición para su uso por los usuarios finales.}
}

\newglossaryentry{Kanban}
{
	name= Kanban,
	description={Método de gestión visual del flujo de trabajo que se caracteriza por el uso de tableros visuales y tarjetas para representar las tareas y su progreso a través de diferentes etapas, promoviendo la transparencia y la mejora continua.}
}

\newglossaryentry{Metodología Ágil}
{
	name= Metodología ágil,
	description={Enfoque iterativo e incremental para la gestión de proyectos, especialmente popular en el desarrollo de software, que se centra en la flexibilidad, la colaboración y la entrega continua de valor. En lugar de adherirse rígidamente a un plan predefinido, agile abraza el cambio y la adaptación a medida que avanza el proyecto.}
}

% Ecuaciones
\usepackage{amsmath}

% Rutas de fichero / paquete
\newcommand{\ruta}[1]{{\sffamily #1}}

% Párrafos
\nonzeroparskip

% Huérfanas y viudas
\widowpenalty100000
\clubpenalty100000

% Imágenes

% Comando para insertar una imagen en un lugar concreto.
% Los parámetros son:
% 1 --> Ruta absoluta/relativa de la figura
% 2 --> Texto a pie de figura
% 3 --> Tamaño en tanto por uno relativo al ancho de página
\usepackage{graphicx}
\newcommand{\imagen}[3]{
	\begin{figure}[!h]
		\centering
		\includegraphics[width=#3\textwidth]{#1}
		\caption{#2}\label{fig:#1}
	\end{figure}
	\FloatBarrier
}

% Comando para insertar una imagen sin posición.
% Los parámetros son:
% 1 --> Ruta absoluta/relativa de la figura
% 2 --> Texto a pie de figura
% 3 --> Tamaño en tanto por uno relativo al ancho de página
\newcommand{\imagenflotante}[3]{
	\begin{figure}
		\centering
		\includegraphics[width=#3\textwidth]{#1}
		\caption{#2}\label{fig:#1}
	\end{figure}
}

% El comando \figura nos permite insertar figuras comodamente, y utilizando
% siempre el mismo formato. Los parametros son:
% 1 --> Porcentaje del ancho de página que ocupará la figura (de 0 a 1)
% 2 --> Fichero de la imagen
% 3 --> Texto a pie de imagen
% 4 --> Etiqueta (label) para referencias
% 5 --> Opciones que queramos pasarle al \includegraphics
% 6 --> Opciones de posicionamiento a pasarle a \begin{figure}
\newcommand{\figuraConPosicion}[6]{%
  \setlength{\anchoFloat}{#1\textwidth}%
  \addtolength{\anchoFloat}{-4\fboxsep}%
  \setlength{\anchoFigura}{\anchoFloat}%
  \begin{figure}[#6]
    \begin{center}%
      \Ovalbox{%
        \begin{minipage}{\anchoFloat}%
          \begin{center}%
            \includegraphics[width=\anchoFigura,#5]{#2}%
            \caption{#3}%
            \label{#4}%
          \end{center}%
        \end{minipage}
      }%
    \end{center}%
  \end{figure}%
}

%
% Comando para incluir imágenes en formato apaisado (sin marco).
\newcommand{\figuraApaisadaSinMarco}[5]{%
  \begin{figure}%
    \begin{center}%
    \includegraphics[angle=90,height=#1\textheight,#5]{#2}%
    \caption{#3}%
    \label{#4}%
    \end{center}%
  \end{figure}%
}
% Para las tablas
\newcommand{\otoprule}{\midrule [\heavyrulewidth]}
%
% Nuevo comando para tablas pequeñas (menos de una página).
\newcommand{\tablaSmall}[5]{%
 \begin{table}
  \begin{center}
   \rowcolors {2}{gray!35}{}
   \begin{tabular}{#2}
    \toprule
    #4
    \otoprule
    #5
    \bottomrule
   \end{tabular}
   \caption{#1}
   \label{tabla:#3}
  \end{center}
 \end{table}
}

%
% Nuevo comando para tablas pequeñas (menos de una página).
\newcommand{\tablaSmallSinColores}[5]{%
 \begin{table}[H]
  \begin{center}
   \begin{tabular}{#2}
    \toprule
    #4
    \otoprule
    #5
    \bottomrule
   \end{tabular}
   \caption{#1}
   \label{tabla:#3}
  \end{center}
 \end{table}
}

\newcommand{\tablaApaisadaSmall}[5]{%
\begin{landscape}
  \begin{table}
   \begin{center}
    \rowcolors {2}{gray!35}{}
    \begin{tabular}{#2}
     \toprule
     #4
     \otoprule
     #5
     \bottomrule
    \end{tabular}
    \caption{#1}
    \label{tabla:#3}
   \end{center}
  \end{table}
\end{landscape}
}

%
% Nuevo comando para tablas grandes con cabecera y filas alternas coloreadas en gris.
\newcommand{\tabla}[6]{%
  \begin{center}
    \tablefirsthead{
      \toprule
      #5
      \otoprule
    }
    \tablehead{
      \multicolumn{#3}{l}{\small\sl continúa desde la página anterior}\\
      \toprule
      #5
      \otoprule
    }
    \tabletail{
      \hline
      \multicolumn{#3}{r}{\small\sl continúa en la página siguiente}\\
    }
    \tablelasttail{
      \hline
    }
    \bottomcaption{#1}
    \rowcolors {2}{gray!35}{}
    \begin{xtabular}{#2}
      #6
      \bottomrule
    \end{xtabular}
    \label{tabla:#4}
  \end{center}
}

%
% Nuevo comando para tablas grandes con cabecera.
\newcommand{\tablaSinColores}[6]{%
  \begin{center}
    \tablefirsthead{
      \toprule
      #5
      \otoprule
    }
    \tablehead{
      \multicolumn{#3}{l}{\small\sl continúa desde la página anterior}\\
      \toprule
      #5
      \otoprule
    }
    \tabletail{
      \hline
      \multicolumn{#3}{r}{\small\sl continúa en la página siguiente}\\
    }
    \tablelasttail{
      \hline
    }
    \bottomcaption{#1}
    \begin{xtabular}{#2}
      #6
      \bottomrule
    \end{xtabular}
    \label{tabla:#4}
  \end{center}
}

%
% Nuevo comando para tablas grandes sin cabecera.
\newcommand{\tablaSinCabecera}[5]{%
  \begin{center}
    \tablefirsthead{
      \toprule
    }
    \tablehead{
      \multicolumn{#3}{l}{\small\sl continúa desde la página anterior}\\
      \hline
    }
    \tabletail{
      \hline
      \multicolumn{#3}{r}{\small\sl continúa en la página siguiente}\\
    }
    \tablelasttail{
      \hline
    }
    \bottomcaption{#1}
  \begin{xtabular}{#2}
    #5
   \bottomrule
  \end{xtabular}
  \label{tabla:#4}
  \end{center}
}



\definecolor{cgoLight}{HTML}{EEEEEE}
\definecolor{cgoExtralight}{HTML}{FFFFFF}

%
% Nuevo comando para tablas grandes sin cabecera.
\newcommand{\tablaSinCabeceraConBandas}[5]{%
  \begin{center}
    \tablefirsthead{
      \toprule
    }
    \tablehead{
      \multicolumn{#3}{l}{\small\sl continúa desde la página anterior}\\
      \hline
    }
    \tabletail{
      \hline
      \multicolumn{#3}{r}{\small\sl continúa en la página siguiente}\\
    }
    \tablelasttail{
      \hline
    }
    \bottomcaption{#1}
    \rowcolors[]{1}{cgoExtralight}{cgoLight}

  \begin{xtabular}{#2}
    #5
   \bottomrule
  \end{xtabular}
  \label{tabla:#4}
  \end{center}
}



\graphicspath{ {./img/} }

% Capítulos
\chapterstyle{bianchi}
\newcommand{\capitulo}[2]{
	\setcounter{chapter}{#1}
	\setcounter{section}{0}
	\setcounter{figure}{0}
	\setcounter{table}{0}
	\chapter*{\thechapter.\enskip #2}
	\addcontentsline{toc}{chapter}{\thechapter.\enskip #2}
	\markboth{#2}{#2}
}

% Apéndices
\renewcommand{\appendixname}{Apéndice}
\renewcommand*\cftappendixname{\appendixname}

\newcommand{\apendice}[1]{
	%\renewcommand{\thechapter}{A}
	\chapter{#1}
}

\renewcommand*\cftappendixname{\appendixname\ }

% Formato de portada
\makeatletter
\usepackage{xcolor}
\newcommand{\tutor}[1]{\def\@tutor{#1}}
\newcommand{\course}[1]{\def\@course{#1}}
\definecolor{cpardoBox}{HTML}{E6E6FF}
\def\maketitle{
  \null
  \thispagestyle{empty}
  % Cabecera ----------------
\noindent\includegraphics[width=\textwidth]{cabecera}\vspace{1cm}%
  \vfill
  
  % Título proyecto y escudo informática ----------------
  \colorbox{cpardoBox}{%
    \begin{minipage}{.8\textwidth}
      \vspace{.5cm}\Large
      \begin{center}
      \textbf{TFG del Grado en Ingeniería Informática}\vspace{.6cm}\\
      \textbf{\LARGE\@title{}}
      \end{center}
      \vspace{.2cm}
    \end{minipage}

  }%
  \hfill\begin{minipage}{.20\textwidth}
    \includegraphics[width=\textwidth]{escudoInfor}
  \end{minipage}
  \vfill
  
  % Datos de alumno, curso y tutores ------------------
  \begin{center}%
  {%
    \noindent\LARGE
    Presentado por \@author{}\\ 
    en Universidad de Burgos --- \@date{}\\
    Tutor: \@tutor{}\\
  }%
  \end{center}%
  \null
  \cleardoublepage
  }
\makeatother

\newcommand{\nombre}{Francisco J. Arroyo Redondo}


% Datos de portada
\title{\fontsize{18pt}{22pt}\selectfont Gestión de Contratos GeNomIn\\
	\fontsize{16pt}{18pt}\selectfont Aplicación APEX sobre Oracle, para la gestión de contratos del Personal contratado con cargo a proyectos de Investigación}
\author{\nombre}
\tutor{Pedro Renedo Fernández}
\date{Septiembre 2025}

\begin{document}

\maketitle


%\newpage\null\thispagestyle{empty}\newpage


%%%%%%%%%%%%%%%%%%%%%%%%%%%%%%%%%%%%%%%%%%%%%%%%%%%%%%%%%%%%%%%%%%%%%%%%%%%%%%%%%%%%%%%%
\thispagestyle{empty}


\noindent\includegraphics[width=\textwidth]{cabecera}\vspace{1cm}

\noindent D. Pedro Renedo Fernández, profesor del departamento de Ingeniería Informática, área de Lenguajes y Sistemas Informáticos.

\noindent Expone:

\noindent Que el alumno D. \nombre, con DNI 13130859, ha realizado el Trabajo final de Grado en Ingeniería Informática. 

\noindent Y que dicho trabajo ha sido realizado por el alumno bajo la dirección del que suscribe, en virtud de lo cual se autoriza su presentación y defensa.

\begin{center} %\large
En Burgos, {\large \today}
\end{center}

\vfill\vfill\vfill

\begin{center}
  Vº. Bº. del Tutor:\\[2cm]
  D. Pedro Renedo Fernández
  \end{center}


\newpage\null\thispagestyle{empty}\newpage




\frontmatter

% Abstract en castellano
\renewcommand*\abstractname{Resumen}
\begin{abstract}
Gestión de Contratos GeNomIn, es una aplicación web desarrollada unen \acrfull{APEX} 24.2 para el control y gestión de los pagos de contratos de  Investigación que actualmente se gestionan a través de una hoja de cálculo. Estos contratos se vinculan a convocatorias y éstas a proyectos de investigación, pretendiendo ificar todas estas actividades en un único aplicativo más sencillo de gestionar por el personal del Servicio y que permita a través de un informe de nómina mensual cotejar los datos remitidos por retribuciones para ese mismo mes, con los calculados por el Servicio.

\end{abstract}

\renewcommand*\abstractname{Descriptores}
\begin{abstract}
Gestión de Pagos, Contratos, Convocatorias, Investigación, Oracle, Apex, GeNomIn
\end{abstract}

\clearpage

% Abstract en inglés
\renewcommand*\abstractname{Abstract}
\begin{abstract}
GeNomIn Contract Management is a web application developed in APEX 24.2 for the control and management of research contract payments, currently managed through a spreadsheet. These contracts are linked to calls for proposals, and these are linked to research projects. The aim is to integrate all these activities into a single application that is easier for Service staff to manage and that allows, through a monthly payroll report, to compare the data submitted for remuneration for that same month with the data calculated by the Service.
\end{abstract}

\renewcommand*\abstractname{Keywords}
\begin{abstract}
Payment Management, Contracts, Calls, Research, Oracle, Apex, GeNomIn
\end{abstract}



\clearpage

% Indices
\tableofcontents

\clearpage

\listoffigures

\clearpage

\listoftables
\clearpage

\lstlistoflistings
\clearpage

\mainmatter
\capitulo{1}{Introducción}

El Servicio de Gestión de la Investigación realiza diversas tareas en el ámbito universitario que van desde la tramitación de facturas, ayudas a la investigación, doctorado, gestión de proyectos, becas y contratos de investigación, pudiendo considerarse en sí misma una "\textbf{mini-gestora}" universitaria.

En este contexto cobran una gran trascendencia los contratos de investigación, ya que generan un gran volumen de trabajo e información. Así, para a realización de un contrato de investigación asociado a un proyecto, entran en liza diferentes ámbitos del servicio que pueden llegar ser inmanejables tal y como se realiza actualmente.

Se manejan aproximadamente 400 contratos de investigación, contabilizando su detalle a través de una hoja de cálculo, que contiene los principales datos, junto con las retribuciones mensuales.

Se pretende con este trabajo, ya no sólo facilitar el trabajo diario, si no evitar errores que se producen fácilmente con el manejo tan voluminoso de hojas de cálculo.
\capitulo{2}{Objetivos del proyecto}

\section{Objetivos funcionales}

Estos objetivos se centran en las funcionalidades y características que debe tener la aplicación \textit{GeNomIn} para satisfacer las necesidades y expectativas de los usuarios. A continuación se detallan los objetivos funcionales del proyecto:

\begin{itemize}
	\item \textbf{Sustitución de la hoja de cálculo Excel}: Uno de los primeros y fundamentales pasos es la conversión de el libro de Excel tablas para la \acrfull{ODB} sin pérdida de datos a través del proceso de \gls{Normalización}.
	En este proceso se irán generado \acrfull{FN} 1, 2 y 3 que nos permitirán; la 1ª eliminar grupos repetidos de datos, y que cada columna tenga un solo valor, la 2ª eliminará redundancias en claves primarias y la 3ª eliminará dependencias transitivas.
	\item \textbf{Creación de las tablas en la \acrfull{ODB}}: Una vez normalizada el libro de Excel, se crearán las correspondientes tablas en la \acrshort{ODB}, configurando correctamente las relaciones entre las mismas \gls{MER}, de vital importancia, ya que una mala definición de las relaciones entre las distintas tablas puede dar al traste con todo el proyecto.
	\item \textbf{Creación de la aplicación sobre \acrfull{APEX}} 24.2: Uno de los principales retos es la creación de dicha aplicación, su correcto funcionamiento y  funcionalidad, que exige el aprendizaje de esta potente  plataforma de desarrollo de aplicaciones web de bajo código que se ejecuta dentro de una base de datos Oracle. Esta versión introduce nuevas funcionalidades, mejoras y actualizaciones, especialmente en el área de la inteligencia artificial generativa y la experiencia del desarrollador.
		
\end{itemize}

\section{Objetivos no funcionales}

Los objetivos no funcionales se refieren a los desafíos y metas que se deben abordar para desarrollar el software. Estos objetivos abarcan aspectos como la arquitectura del sistema, las tecnologías a utilizar y las metodologías de desarrollo. A continuación se detallan los objetivos no funcionales del proyecto:

\begin{itemize}
	\item \textbf{Gestión de las diversas herramientas para la creación de la aplicación GeNomIn}: 
	Durante la asignatura de \acrfull{AGBD}, se han visto someramente las herramientas \acrshort{APEX} 5.1 y Oracle 9. Teniendo el objetivo de desarrollar una aplicación funcional y entendiendo que ha de desplegarse en un servidor \acrshort{HTTPS} con \gls{SLL}, se prefirió como reto personal, el desarrollo sobre \acrfull{APEX} 23.2 sobre Oracle 23ai. También durante la asignatura de \acrfull{VYP}, se han utilizado herramientas de prueba como Selenium y Katalom, se ha querido salir de la zona de confort y utilizar \textbf{textcafe} para comprobar la usabilidad de la aplicación.
	\item \textbf{Usabilidad de \textbf{GeNomIn}}: La aplicación no debe ser sólo útil, como principal objetivo, si no que debe permitir un manejo fácil e intuitivo por el personal que debe manejarla. Como se ha comentado en la introducción la familiaridad del personal con \acrfull{APEX} en las distintas aplicaciones ya usadas (portales Económico e Investigación), favorece una mejor asimilación de este nuevo entorno en comparación con la tediosa hoja de cálculo. 
\end{itemize}

\section{Objetivos personales}

\begin{itemize}
	\item \textbf{Capacidad de desarrollo}: el principal objetivo es comprobar si el esfuerzo realizado durante todos estos años es aplicable a la "vida profesional" y tiene aplicación directa en un caso real. 
	\item \textbf{Reto de aprendizaje}: Lo sencillo hubiera sido realizar este trabajo con herramientas conocidas y se ha querido llevar un paso más allá, utilizando versiones muy superiores a las utilizadas en prácticas (\acrshort{APEX}5.1) y software totalmente distinto para la realización de pruebas.
	También, la realización de esta memoria con \gls{LATEX}, supone un cambio de paradigma sobre el uso de los procesadores texto convencionales como Word.
	\item \textbf{Reto personal}: la realización del \acrshort{TFG} supone el final del camino iniciado hace ya más de 20 años cuando inicié la Ingeniería Técnica Informática y por motivos personales no pude concluir.
\end{itemize}
\capitulo{3}{Conceptos teóricos}
\section{Introducción}
En este capítulo se van a presentar los fundamentos teóricos en los que se basa esta aplicación, cómo nace la necesidad y surge la idea de su desarrollo.

\section{Antecedentes} 
La idea de la realización de este \acrshort{TFG}, surge principalmente de la relación laboral de este estudiante con la Universidad, concretamente en el Servicio de Gestión de la Investigación. Como ya se comentó en la introducción de esta memoria, en este Servicio se realizan diferentes tareas relacionadas
con la investigación y particularmente contratos asociados a proyectos de investigación.

Para la realización de este tipo de contratos los Investigadores, realizan una solicitud de contratación, luego solicitantes con diversas características
requeridas (sobre todo titulación), participan en las mismas y para el que es seleccionado, se realiza un contrato asociado a un proyecto, \textbf{que tiene que
estar en vigor}. 

Hasta aquí, cabría esperar que finalizase este proceso, pero hay que tener en cuenta las disponibilidades presupuestarias de los proyectos y los cambios que pueden surgir durante la vida del contrato.

\section{Problemática}
Una vez generado el contrato, el Servicio de Recursos Humanos genera mensualmente dos nóminas para cada contratado, una de seguridad social y otra de la nómina en sí. Estamos hablando que actualmente hay más de cuatrocientos contratos vinculados a estos proyectos.

Estás nóminas (\textbf{sábanas}), son remitidas al Servicio de Investigación para que de conformidad a la existencia de crédito, la cuantía y validez del proyecto, antes de ordenar efectuar el pago general a contabilidad.\label{img: sabana}
\imagen{sabana}{Sábana de nómina de RRHH}{1
} 

Para ésto, la Sección de Personal Contratado del Servicio de Investigación al realizar un contrato, a punta en una hoja Excel los datos del contratado, el proyecto vinculado, las mensualidades de un año con su seguridad social y observaciones.\label{img: contrato}
\imagen{tablaexcel}{Tabla de nómina}{1}

Cuando recibe las sábanas de nómina, coteja cada referencia del proyecto (\textbf{\textbf{orgánica}}), con los datos que tiene apuntados en su hoja de cálculo y comprueba si hay diferencias, si los datos del proyecto son correctos así como las cuantías totales por orgánica.

Esto proceso puede repetirse varias veces si hay errores. Al finalizar el año se hace una copia de la tabla de Excel del año anterior y se copia en el siguiente, donde se añaden los nuevos contratos.

\section{Proyecto GeNomin}

En este estado de cosas y con la oportunidad de realizar un \acrshort{TFG} de utilidad, se propone transformar la hoja de cálculo en una \acrshort{BDR} en Oracle, que unifique todos los datos que ahora se dispersan en hojas de cálculo anuales y a través de una aplicación en \acrshort{APEX}, permita obtener informes de nómina que faciliten el trabajo de cotejo con las sábanas de retribuciones. Con la particularidad de que la aplicación de Gestión de Proyectos, Convocatorias y Contratos, se desarrolla también sobre ambas plataformas, no supondría un problema de adaptación para los usuarios.

En el proyecto, se podrán gestionar investigadores, proyectos, convocatorias y solicitantes. Asociar un solicitante a una convocatoria y que generar las nóminas correspondientes. Prorrogar y renunciar a los contratos (modificando la nóminas a pagar) Y finalmente generar diversos informes de utilidad.
\capitulo{4}{Técnicas y herramientas}

Otro de los aspectos a tener en cuenta en la realización de cualquier proyecto, ya sea de albañilería o de desarrollo informático, son las herramientas con las que se realiza, ya que de ello dependerá mucho tanto el resultado como la planificación de la ejecución.
En este capítulo se detallarán las herramientas escogidas para el desarrollo de la aplicación \textbf{GeNomIn} y su justificación.

Puesto que la aplicación se desarrolla sobre BD, la primera elección tenía que estar condicionada por qué \acrshort{SGBD} elegir. Como se indica en el capítulo anterior (\ref{sec:SGBD}) entre los más usados comercialmente están; \textbf{MySQL}, \textbf{PostgreSQL}, \textbf{SQL Server} y \textbf{Oracle}. Puesto que este \acrshort{TFG} pretende ser un reflejo de las competencias adquiridas durante el Grado y dado que en la asignatura de \acrshort{SGBD}, se usa también Oracle 9 parece una opción elegible. También se estudió el sistema empleado por la Universidad de Burgos para el almacenamiento de sus BD, siendo también Oracle, con lo cual fue la opción elegida.
Siguiendo este mismo criterio de continuidad, para el desarrollo de la aplicación se escogió \acrshort{APEX} 5.1, ya que era la plataforma utilizada en el Grado.

Expuesto el proyecto al Tutor y puesto que los objetivos de los \acrshort{TFG} deben ser también una mejora en las competencias adquiridas, se optó por una versión más moderna que permitiera el despliegue completo en cloud, por lo que definitivamente se escogieron como principales herramientas Oracle23ai Free y \acrshort{APEX}2402.

Con esta decisión también fue necesaria la configuración de \acrshort{ORDS} y \acrshort{HTTPS}, como se detallará más adelante.

\section{Oracle23ai}
\imagen{Oracle23ai}{Esquema de Oracle23ai}{.6}

\textbf{Oracle23ai Free} es una de las plataformas más usadas en el panorama comercial y que sigue ofreciendo su \acrshort{IDE}, SQL Developer y varias funcionalidades de forma gratuita para desarrolladores.
Entre las principales características se encuentran:
\begin{itemize}
	\item \textbf{AI Vector Search:} (búsqueda vectorial de AI) es una colección de funciones que incluye un nuevo tipo de datos vectoriales, índices vectoriales y operadores SQL de búsqueda vectorial que permiten a \acrfull{ODB} almacenar el contenido semántico de documentos, imágenes y otros datos no estructurados como vectores y utilizarlos para ejecutar consultas de similitud rápidas.
	\item \textbf{JSON Relational Duality Views:} las vistas de dualidad relacional de JSON unifican los modelos de datos relacionales y documentales para ofrecer lo mejor de ambos mundos.
	\item \textbf{Gráficos de propiedades operativas:} Ofrece soporte nativo para estructuras de datos de gráficos de propiedades y consultas de gráficos.
	\item \textbf{SQL Firewall:} función de seguridad de base de datos integrada en el núcleo de \acrfull{ODB} que inspecciona todas las conexiones entrantes a la base de datos y sentencias SQL, y permite/registra/bloquea actividades no autorizadas de acuerdo con políticas específicas del usuario de base de datos.
	\item \textbf{True Cache:} esta solución simplifica el almacenamiento en caché en \acrfull{ODB}.
	\item \textbf{Mejoras de SQL:} incluye nuevas funciones, como dominios de uso de aplicaciones, que permiten a los desarrolladores definir las columnas que representan.
	\item \textbf{Mejoras en la escalabilidad y disponibilidad}
\end{itemize}

El proceso de desarrollo fue instalar localmente ~\cite{DattaInstallingOracleDatabasea} para una vez finalizado el producto desplegarlo en la nube.

\section{Apex 24.02} \label{sec: apex}
Para el desarrollo de \textbf{GeNomIn} se ha utilizado, como se ha comentado anteriormente \acrshort{APEX} 2402, instalado inicialmente también de forma local ~\cite{JenningsInstallingConfiguringAPEX}
\imagen{apex}{Home Apex TFGUBU-GeNomIn}{.6}
Oracle Apex 24 es una plataforma \textbf{low-code} para desarrollos empresariales más utilizada que permite crear aplicaciones móviles y web escalables y seguras, pudiéndose implementar en la nube o localmente de forma gratuita.

Sus principales características son:
\begin{itemize}
	\item \textbf{IA Generativa}: ofrece capacidades de IA generativa mejoradas, incluyendo un asistente de IA más potente y la posibilidad de configurar datos RAG (Retrieval-Augmented Generation)
	\item \textbf{Database Object Dependencies}: permite escanear y revisar los Objetos de Base de datos que están siendo referenciados en nuestra aplicación.
	\item \textbf{Fuentes de datos REST mejoradas}: amplía su capacidad para integrarse con APIs externas y servicios web de forma más flexible y potente.
	\item \textbf{Soporte para fuentes JSON}: permite trabajar con datos en formato JSON sin necesidad de almacenarlos en tablas tradicionales.
	
\end{itemize}

\section{\acrshort{ORDS}}
Oracle \acrshort{REST} Data Service Oracle REST Data Services sirve de puente entre \acrshort{HTTPS} y tu Oracle Database. ORDS, una aplicación Java de nivel medio, proporciona una API REST de gestión de bases de datos, SQL Developer Web, una puerta de enlace PL/SQL, SODA para REST y la capacidad de publicar servicios web RESTful para interactuar con los datos y los procedimientos almacenados en su Oracle Database.

Puesto que era necesario el acceso \acrshort{HTTPS}, fue necesaria la generación e instalación de \acrshort{SSL}, para ello se generó un certificado con OpenSSL y su posterior con configuración, como se muestra en la imagen.
\imagen{ords_cfg}{Configuración ORDS-HTTPS}{.5}

\section{TestCafe}
TestCafe es un \acrshort{IDE} multiplataforma para pruebas web integrales que no requiere WebDriver ni otras herramientas. TestCafé Studio funciona en Windows, macOS y Linux, y puede ejecutar pruebas en cualquier navegador de escritorio o móvil. ~\cite{TestCodeGuide}
Principalmente se escogió este \acrshort{IDE}, como en el caso de la versión mejorada de \acrshort{APEX}, por ofrecer un desarrollo en entornos no manejados durante el Grado, tales como \textbf{Selenium}. 
Las principales características que ofrece son:
\begin{itemize}
	\item \textbf{Grabación de pruebas y captura de imágenes}: TestCafe permite la grabación y toma de imágenes de las pruebas efectuadas, pero en la práctica, se limita a unos pocos frames inicales.
	\item \textbf{Habilitado para distintos navegadores}: las pruebas pueden ser ejecutadas en diversos navegadores, chrome, edge, firefox, safari, etc..
	\item \textbf{Conjunto completo de Assertions}: en cualquier momento se puede comprobar el estado de un elemento, contenido, posición, valores de retorno y demás propiedades.
	\item \textbf{Espera automática}: está diseñado para la web asíncrona moderna. Identifica correctamente eventos como la carga de páginas, la representación de elementos y las solicitudes XHR, y espera a que se resuelvan antes de continuar con la prueba.
	\item \textbf{Selectores automáticos}: Las acciones y aserciones de prueba utilizan consultas de selector para identificar elementos por su clase, atributos, ID u otras propiedades. Cuando se registra una prueba, se genera consultas de selector automáticamente.
	\item \textbf{Informes completos}:  genera informes de pruebas completos que incluyen resúmenes de las ejecuciones, así como información detallada sobre cada prueba. En nuestro desarrollo se han utilizado los informes \textbf{Allure}, ver imagen(\ref{img: kanban})
\end{itemize}

En la código siguiente vemos un ejemplo de test 25 que comprueba si se muestra el informe entre fechas y en la imagen su ejecución:

\begin{lstlisting}[language=JavaScript, caption={Codigo en js para prueba en TestCafe}]
	import { Selector } from 'testcafe';
	import {NOMINAMENU, BUTTON, USERNAME, PASSW, TOGICON4, BT_CONNOM, VTOMENU } from './constanst.js';
	
	fixture`Test Suite-25`.page("https://192.168.2.61:8443/apex/f?p=100:LOGIN_DESKTOP:12651011480748:::::")
	.meta({TEST_RUN:'Informe de contratos entre fechas',FEATURE: 'Informes', STORY: 'US25-Zube #23'});
	
	test.meta({SEVERITY:'critical', ISSUE_URL: 'https://github.com/far0010/TFGUBU-Fran_Arroyo/issues/16',
		STORY: 'US25-Zube #23'})
	('US25-Zube #23: Informe de contratos entre fechas', async t => {
		const INI = '01-ENE-2025'; 
		const FIN = '31-DIC-2025';
		// Espera a que los campos esten disponibles
		await t.expect(USERNAME.exists).ok({ timeout: 5000 });
		await t.expect(PASSW.exists).ok({ timeout: 5000 });
		
		// Interaccion con los elementos
		await t
		.typeText(USERNAME, 'user01')
		.typeText(PASSW, 'user01')
		
		await BUTTON();;
		await t
		.click(TOGICON4)  // Hace clic en el icono de despliegue
		// Esperar a que aparezca Nomina mes
		await t.expect(VTOMENU.exists).ok();
		
		// Hace clic en el enlace Nomina mes
		await t.click(VTOMENU);
		// introducir fechas inicio - fin
		const BT_VTOS = Selector('#B18949343265529115')
		await t
		.typeText(Selector('input[name="P17_FDESDE"]'), INI)
		.typeText(Selector('input[name="P17_FHASTA"]'), FIN)
		.click(BT_VTOS);
		
		// comprobamos que ofrece las filas deseadas 3
		const filasDatos = Selector('#R18949894423529120_data_panel tbody tr')
		.filter(node => node.querySelectorAll('td').length > 0);
		await t
		.expect(filasDatos.count)
		.eql(3, 'La tabla no muestra exactamente 3 filas de datos');        
	});
\end{lstlisting}

\imagen{test25}{Resultado del Test 25:comprueba el informe entre fechas}{.6}

Una vez realizado el desarrollo y completadas las pruebas establecidas en los \gls{CasodeUso}, se realiza el despliegue. Así, siguiendo la misma operativa en el desarrollo del \acrshort{TFG}, se opta por el despliegue en herramientas gratuitas, en este caso Oracle sigue ofreciendo su versión \textbf{cloud free tier}, para desarrolladores \href{https://www.oracle.com/es/cloud/free/}{\acrfull{OCI}}).

\section{Oracle Cloud Free Tier}\label{sec: OFT}
Como se ha indicado anteriormente \acrshort{OCI}, ofrece una serie de servicios "\textit{gratuitos}" para desarrolladores en su nube. Tras el proceso de registro, que no es nada sencillo, ya que es preciso aportar una tarjeta física en vigor (con un cobro de 0,73€) y los datos son revisados. Para este trabajo especialmente se accede a:
\begin{itemize}
	\item \textbf{Autonomus Data Base}: Base de datos autónoma sin costo que se gestiona automáticamente (aprovisionamiento, seguridad, disponibilidad, rendimiento, etc.)
	\item \textbf{AMD Compute Instance}: 2 \acrshort{VM} basadas en AMD, con 1/8 de OCPU y 1 GB de memoria cada una.
	\item \textbf{\acrshort{APEX}}: plataforma para el desarrollo de aplicaciones low-code 
\end{itemize}

\imagen{OracleCloud}{Base de Datos Autónoma en Oracle Cloud}{.6} \label{img: OCloud}

Una vez creada la instancia de la base de datos autónoma, es preciso crear la instancia de \acrshort{APEX} e iniciar el servicio y así poder proceder a la importación tanto de la BD como la aplicación generada localmente.
Para el proceso de exportación se utilizó el asistente de SqlDeveloper e importando los datos en Cloud a través de las acciones de la BD, sql en su hoja de trabajo. Se intentó realizar el proceso a través de pump de datos (que realiza una copia especular), pero principalmente por compatibilidades de configuración (zona horaria) resultó imposible.

\section{Apex Office Print}
Uno de los objetivos principales de la aplicación, además de la gestión de los contratos, es ofrecer informe detallado de las nóminas que se pagan en un mes determinado. Queriendo ofrecer una personalización de los mismo se utiliza el servicio \acrshort{AOP}, que permite generar informes Office y PDF, pudiendo personalizar éstos ~\cite{OverviewAPEXOffice}.
Para ello es necesario darse de alta en el servicio y generar una plantilla, en la que se detallen los campos del informe a imprimir:
\\
\imagen{infAOP}{Plantilla de informes-Informe AOP}{.6} \label{img:plantilla}

\section{GitHub}
Es una plataforma para alojamiento y gestión de código fuente que se basa en el sistema de control de versiones Git, siendo una herramienta fundamental en el desarrollo de proyectos ~\cite{ChaconProGitTodo}
Para este proyecto se ha utilizado junto con la versión de escritorio \textbf{GitHub Desktop}, haciendo commits tanto del nuevo software como de las actualizaciones y subiéndolas al repositorio (push).
\imagen{github}{GitHub del Proyecto GeNomIn}{.6}
\subsubsection{Control de Versiones}:
las diferentes versiones del desarrollo realizado, están disponibles a través de los commits. Principalmente se han actualizado código \textbf{SQL}, \textbf{PL/SQL} y \textbf{JavaScript}, además de la diversa documentación generada. 
Es posible la solicitud de cambios de código al resto de colaboradores (pull-requests), aunque en el ámbito de este desarrollo individual no tiene mucho sentido.
\subsubsection{Integración de Herramientas}
Como se puede observar en la imagen anterior, \textbf{github}, permite la integración de diversas herramientas que permiten un mayor control del proyecto, tales como \textbf{SonarCloud}, para verificar el código a través de \textit{Actions}, \textbf{Zube} e informes de \textbf{Allure}
\subsubsection{Documentación}
El repositorio permite la organización de la documentación de forma estructurada, así como la creación de \textit{Wikis}, para facilitar la comprensión del proyecto.
\subsection{\gls{Release}}
Una de las partes importantes de un proyecto, es la generación de versiones que indican el progreso del mismo. GitHub facilita esta tarea con su funcionalidad de \textit{releases} que permiten empaquetar y distribuir las distintas versiones del producto de forma organizada y documentada.
\imagen{release}{Vista de un release}{.6}

\href{https://github.com/far0010/TFGUBU-Fran_Arroyo}{Proyecto GeNomIn en GitHub}

\section{Zube}
Zube es una plataforma para gestión de proyectos colaborativos. Se basa en tableros \gls{Kanban} que facilitan la planificación, la gestión de \textbf{sprints}, \gls{Milestone} y demás herramientas gráficas que permiten un análisis del trabajo a realizar y realizado. Ver imagen(\ref{img: kanban})

Zube está perfectamente integrado con \textbf{GitHub} permitiendo una sincronización de tareas y código fundamentales para la planificación del proyecto.

\section{TextStudio}

TextStudio es una herramienta especializada en la edición de textos \textbf{\gls{LATEX}} que facilita la redacción de trabajos académicos y científicos mediante características como el resaltado de sintáxis, corrección ortográfica y semántica en tiempo real y fundamentalmente, la \textit{compilación automática}, la cual permite previsaulizar el documento mientras se trabaja. 
Es cierto que al estar acostumbrados a los productos office, inicialmente resulta "inquietante", pero los resultados son más profesionales ~\cite{BibliotecaComplutenseLaTeXTuTFG2024}.

\capitulo{5}{Aspectos relevantes del desarrollo del proyecto}
En este capítulo se detallan los aspectos que han influido en el desarrollo de la aplicación \textbf{GeNomIn}, tanto técnicas, económicas y de oportunidad.

\section{Oportunidad de negocio}
Cuando se aborda un \acrshort{TFG}, según mi opinión, debe poder trasladarse al ámbito profesional y así desarrollar las competencias adquiridas durante estos años de aprendizaje en el Grado.
En el entorno laboral en el que desarrollo mi trabajo, existía la necesidad de un control más férreo y menos tedioso de los pagos de nómina a través de proyectos de Investigación que se estaban realizando a través de hojas de cálculo. En este sentido, cabe explicar que no es que se pague mal, si no en el proceso de comprobación utilizado, en  lo relativo a los contratos realizados con proyectos de investigación.
Mensualmente se emite un listado de los pagos que corresponden a estos proyectos y se cotejan contra una hoja de Excel, y con más de 400 contratos, es evidente que resulta inviable. (ver imagen \ref{img: contrato})
Así pues, se planteó la propuesta de crear una aplicación, que no sólo ofreciera un tipo de listado más uniforme, si no que permitiera un control más adecuado de las personas contratadas, las convocatorias y sus responsables.
En este sentido, hay que indicar, que la plataforma que actualmente utiliza la Universidad, ofrece la gestión de convocatorias y contratos, pero no es posible el control de pagos de estos contratos con los módulos instalados.

\section{¿Por qué APEX 2402?}
Uno de los aspectos más importantes y que influyen en el desarrollo de una aplicación es la herramienta que se utiliza. En este sentido y puesto que el entorno de trabajo se ha desarrollado con \acrshort{APEX}, parecía bastante adecuado, tratando así de que la transición de la hoja de cálculo a la aplicación sea lo menos costosa para los usuarios.

Como vimos en la sección dedicada a esta herramienta \ref{sec: apex},  \acrshort{APEX} ofrece muchas características técnicas que ya de por sí solas la hacen elegible, pero veamos qué facilita para el desarrollo y la gestión.
\subsection{Integración con \acrfull{ODB}}
\acrshort{APEX} aprovecha la potencia y eficiencia de \acrshort{ODB} así como sus características avanzadas así como integración nativa con \acrshort{SQL} y \acrshort{PL/SQL} desarrollado por Oracle ~\cite{OracleAPEXSQL}
\subsubsection{Gestión de Usuarios}
\imagen{apex_acl}{Gestión de usuarios de APEX}{1}
Como vemos en la imagen anterior \acrshort{APEX}, permite un control total de los usuarios a través de:
\begin{itemize}
	\item \textbf{\acrfull{ACL}}: permiten asignar roles asuarios específicos.
	\item \textbf{Esquemas de autorización}: que se aplican a páginas o componentes individuales para restringir acceso.
	\item \textbf{Generación automática de componentes}: al activar este \acrshort{ACL}, se generan automáticamente distintos componentes; administración de usuarios, región de control de acceso, roles, compilaciones condicionales, etc. ~\cite{OracleAPEXApp}, \cite{OracleAPEXAdministration}
\end{itemize}
\section{Creación de las Páginas}
\acrshort{APEX} combina el desarrollo visual con la lógica declarativa y programación libre, para el desarrollo de aplicaciones web a través del \textbf{diseñador de páginas} (entorno visual para definir la estructura y componentes) y del \textbf{asistente de creación} (para generar formularios, informes interactivos, dashboards, etc.)
\imagen{apex_pag}{Vista de las páginas APEX}{1}
\section{Programación}
\acrshort{APEX} permite la programación de procesos y acciones dinámicas en \acrshort{PL/SQL}, \acrshort{JS} sin complejidad de código.
Además se pueden establecer condiciones y validaciones para mostrar u ocultar elementos según reglas establecidas.
\imagen{apex_act}{Acciones dinámicas APEX y su código js}{1}
En la imagen superior podemos ver un grupo de acciones dinámicas asociadas a la pag 12 del proyecto: \textbf{Nuevo contratado}, y se muestra el código de la que verifica si la fecha de inicio del contrato, es posterior a la fecha de inicio del proyecto asociado.
\section{Despliegue en Oracle Cloud Free Tier}
Una de las características principales por las que se eligió este modelo de desarrollo en el binomio \acrshort{APEX} \ \acrshort{ODB}, fue la gratuidad y facilidad de los servicios. Aunque como ya se explicó en el capítulo 4 (ver \ref{sec: OFT}), el alta en Cloud, no es sencillo, se prefirió a otras opciones como \textbf{kubernetes}, con una curva de aprendizaje y configuración mayores o \textbf{virtualbox} que hubiera valido para un entorno local, exportando la imagen y confiando en la configuración del usuario de destino.

Así, para este tipo de desarrollos más \textit{experimentales}´ se recomienda su uso, ya que además permite interactuar con entorno profesional, de bajo coste, alta disponibilidad, escalable (previo pago) y seguridad integrada, todo ello ofrecido por Oracle Cloud. Ver imagen(\ref{img: OCloud})



\capitulo{6}{Trabajos relacionados}
Para la realización de este trabajo se hicieron diversas consultas a través de buscadores, de aplicaciones similares, que hicieran esta función de aunar los proyectos de investigación, sus convocatorias y sus contratos emitiendo informes personalizados.
Así, existen en el mercado aplicaciones como \href{https://www.pandadoc.com/}{PandaDoc} o \href{https://www.zoho.com/es-xl/contracts/?lb=es-xl}{Zoho Contracts} que gestionan contratos, pagos, etc..., pero no están vinculados a convocatorias de proyectos de investigación.

En este sentido y ya conocido,  \href{https://www.universitasxxi.com/}{Universitas XXI}, sí ofrece varios módulos dedicados a la gestión de RRHH, proyectos de Investigación y gestión económica, con emisión de varios informes independientes, pero que no se adaptan por completo a los requerimientos de los usuarios.
Veamos en esta tabla comparativa entre la aplicación \textbf{GeNomIn} y los módulos de \textbf{\acrshort{UXXI}}:

\begin{table}[h!]
	\centering
	\footnotesize
	\begin{tabular}{|l|l|l|l|l|}
		\hline
		\rowcolor{gray!20}
		\multicolumn{1}{c}{\textbf{Funcionalidad}\rule{0pt}{25pt}} & \multicolumn{1}{c}{\textbf{GeNomIn}} & \multicolumn{1}{c}{\textbf{UXXI-Inv}} &
		\multicolumn{1}{c}{\textbf{UXXI-Eco}} &
		\multicolumn{1}{c}{\textbf{UXXI-RRHH}} 
		\\
		\hline
		Contratos Inv. & Parcial & Sí & No & No\\
		Convocatorias Inv. & Parcial & Sí & No & No\\
		Ges. Investigadores & Parcial & Sí & No & No\\
		Informes nómina & Sí & Sí & General & Contable\\
		Contratos con pagos & Sí & NO & General & Contable\\
		Informes Person. & Sí & Propio & Propio & Propio\\
		Infraestructura & Oracle Cloud & Propio & Propio & Propio\\
		Flexibilidad & Alta & Versiones & Versiones & Versiones\\
		Costo & Free Tier & Licencia & Licencia & Licencia\\
		\hline
	\end{tabular}
	\caption{Comparativa GeNomIn vs Universitas XXI}
	\label{tab:comGenUniv}
\end{table}

Como podemos comprobar la opción comercial tiene muchísimo más desarrollo pero no se personaliza a los objetivos de cada usuario, teniendo que esperar a la liberación de versiones, si es que está incluida esa actualización o solicitar una personalización, que generalmente no es barata.
\capitulo{7}{Conclusiones, valoración personal y trabajo futuro}
En este capítulo se detalla si se han cumplido los Objetivos del proyecto, tanto funcionales, no funcionales y personales, y si se han desarrollado y consolidado las competencias adquiridas durante el Grado.

\section{Cumplimiento de objetivos funcionales}
El principal objetivo de este \acrshort{TFG}, era la transformación de una hoja de cálculo Excel, en la que se contabilizan las nóminas del personal contratado (ver imagen \ref{img: contrato}), en una \acrshort{ODB} y la generación de una aplicación en \acrshort{APEX}, que además de llevar la gestión de las convocatorias y contratos, presentara un informe para el cotejo con la retribución emitida por RRHH (ver imagen \ref{img: sabana}).

Una vez finalizado el desarrollo, el proyecto ha logrado cumplir estos objetivos, desarrollando una aplicación \textbf{Genomin}, que cumple a la perfeción todos los retos planteados.
\imagen{logo}{Logotipo de la aplicación: GeNomIn}{.2}

\section{Cumplimiento de objetivos no funcionales}
Dentro de los objetivos no funcionales se pretendía reforzar las competencias adquiridas durante el Grado, así como explorar nuevas herramientas de desarrollo y ampliar algunas de las conocidas. Todos éstos retos han sidos superados desde la gestión de proyectos con la \gls{Metodología Ágil}, generación de \acrshort{ODB}, \acrshort{SSL}, creación de la aplicación en \acrshort{APEX} 24, realización de test de \acrshort{VYP},  hasta su despliegue en \textbf{Oracle Cloud} y todo ello documentado y apoyado con el repositorio documental \textbf{GitHub}

\section{Valoración Personal}
Antes de empezar este \acrshort{TFG}, tenía la sensación de que era un mero trámite más para la consecución del Grado en Informática, pero lo cierto es que es una de las asignaturas más enriquecedoras. Ya no solo supone un reto en cuanto al desarrollo de código, que hoy en día no lo es tanto con la irrupción de la \acrshort{IA}, si no que supone un esfuerzo organizativo, que requiere de muchas de las competencias adquiridas durante la carrera. Es aquí donde se comprenden muchos conceptos, como la \gls{Normalización}, \gls{MER}, la gestión de proyectos, el uso de \gls{Metodología Ágil}, la \acrfull{VYP} etc., en el entorno de desarrollo de un proyecto.
En mi caso la realización de este proyecto una vez finalizado el curso y completadas el resto de asignaturas, me ha permitido dedicar mucho más tiempo, organización e investigación que de otra manera hubiera resultado muchísimo más complejo.

Con todo esto, la aplicación web \textbf{GeNomIn} desarrollada, desde mi punto de vista, cumple con los objetivos inicialmente planteados y podría ser desplegada perfectamente para su uso en el Servicio para el que ha sido diseñada, puesto que corre sobre \acrfull{ODB}, la BD utilizada por la Universidad y teniendo como plataforma \acrshort{APEX}, utilizado por \acrshort{UXXI}-Investigación (actualmente proveedor de la Universidad de Burgos), lo cual no supondría ningún problema de adaptabilidad por parte de los usuarios.

\section{Trabajo futuro}
El trabajo abordado en este \acrshort{TFG}, nace de la necesidad observada en la gestión y cotejo de pagos al personal investigador contratado, con cargo a proyectos de investigación.
En este sentido se ha abordado la generación de este informe (\textbf{Pago nómina-mes}), otros como la renovación y la renuncia de contratos, e informes adicionales que muestran información sobre vencimientos de contratos en un periodo determinado y de carácter general.

Quedaría por abordar, las \textbf{subidas salariales} y \textbf{bajas médicas} que también tienen influencia en los pagos mensuales, que tal y como está ahora mismo la aplicación, se podrían realizar modificando directamente los pagos en la tabla correspondiente (\textbf{NOMINAS}), pero sin tener una gestión directa.
Seguramente, si se llegase a poner en explotación, sería imprescindible añadir otro tipo de \textbf{informes personalizados}, que surgieran del propio uso de la aplicación y la gestión diaria.

Evidentemente este proyecto simplemente es para uso de forma local, por el  propio personal del Servicio de Investigación que gestiona contratos, ya que la plataforma \acrshort{UXXI} desarrolla sus propios módulos bajo licencia, pero puesto que se ha considerado su necesidad podría ser planteado para desarrollos futuros a la propia empresa.

\section{Reflexión sobre el uso de \acrshort{IA}}
Durante la realización de este trabajo se han realizado diversas consultas a través de \gls{Copilot} sobre todo en "\textbf{momentos de pánico}". En este sentido cabe destacar, que en las versiones libres (la utilizada en este proyecto), son un mero apoyo y consulta, ya que normalmente las respuestas obtenidas ofrecen más una orientación, que una respuesta efectiva y que generalmente pueden llevar al error, sobre todo cuando se realizan, por ejemplo, consultas complejas.
Evidentemente no se pueden poner puertas al campo y el uso de la \acrshort{IA} está cada vez más extendido, sería deseable su inclusión en las asignaturas del Grado, como apoyo al desarrollo y al conocimiento, como en su día fueron, los libros, Internet y ahora esta nueva tecnología. Tendremos que adaptarnos todos.




\printnoidxglossary[type=\acronymtype]
\printnoidxglossary[]
%\printacronyms
%\printglossary
%\printnoidxglossary[type=\acronymtype]
%\printglossary[type=\acronymtype]
%\printglossary



\nocite{*}
\bibliographystyle{plain}
\bibliography{bibliografia}

\end{document}